\section{The dependency graph}
\label{sec:dependency-graph}

\subsection{Two branches and one merge}

The audited chain is
\begin{center}
\begin{tabular}{c}
\(\boxed{\text{TPC--63 provenance}}\)\\[-1mm]
\(\downarrow\)\\[-1mm]
\(\boxed{\text{TPC--64 activation}}\)\\[-1mm]
\(\downarrow\)\\[-1mm]
\(\boxed{\text{TPC--65 represented Gram}}\)\\[-1mm]
\(\downarrow\)\\[-1mm]
\(\boxed{\text{TPC--66 support geometry}}\)\\[-1mm]
\(\downarrow\)\\[-1mm]
\(\boxed{\text{TPC--67 determinant mass}}\)
\end{tabular}
\qquad
\begin{tabular}{c}
\(\boxed{\text{TPC--68 signed native}}\)\\[-1mm]
\(\downarrow\)\\[-1mm]
\(\boxed{\text{TPC--69 first mixed trace}}\)\\[-1mm]
\(\downarrow\)\\[-1mm]
\(\boxed{\text{TPC--70 anchored plane}}\)\\[-1mm]
\(\downarrow\)\\[-1mm]
\(\boxed{\text{TPC--71 abstract variation}}\)\\[2mm]
\(\phantom{\boxed{\text{TPC--71 abstract variation}}}\)
\end{tabular}
\end{center}
The left branch seeks to control the represented map \(A\).  The right
column starts from the missing map \(M\), but after TPC--68 it descends to
the scalar sector \(\tr(ER)\); it does not return to an operator bound for
\(M\).  The full maps merge through
\[
 B=A+M
\]
and a shorted comparison relative to \(A\).

\subsection{The exact reassembly object}

Let
\[
 N_A=\Ker A,
 \qquad
 W_{\rm nuis}=B(N_A)=M(N_A).
\]
The direct represented-normalized Gram is not the final uniform comparison,
because vectors in \(N_A\) can produce physical output through \(M\).  The
exact shorted Gram has the form
\begin{equation}
 G_{\rm sh}=G_{\rm dir}-C_N,
 \qquad
 C_N\succeq0,
 \label{eq:tpc72-shorted-gram}
\end{equation}
following the classical shorted-operator framework
\citep{AndersonTrapp1975} and the TPC--68 specialization
\citep{WangTPC68}.

The native Gram \(I+R+E\) controls an atomic-normalized part of \(M^*M\).
It is not automatically \(G_{\rm dir}\), which is normalized by represented
energy, and it does not control \(C_N\).

\begin{theorem}[Dependency theorem]
\label{thm:dependency}
A uniform endpoint certificate from the audited architecture requires, in
one literal carrier and one quantifier scope:
\begin{enumerate}[label=\textup{(\roman*)},leftmargin=3em]
 \item a positive terminal activation arrow;
 \item a represented lower frame for \(A\);
 \item a missing-mass comparison;
 \item an upper bound for the complete missing-native Gram;
 \item a negligible or explicitly charged nuisance penalty \(C_N\);
 \item boundary and tail closure; and
 \item a total physical loss strictly below \(1/400\).
\end{enumerate}
No proper subset of these requirements is supplied as the full endpoint
theorem by TPC--63--TPC--71.
\end{theorem}

\begin{proof}
TPC--64 shows that a terminal-inactive row kills the uniform first arrow.
TPC--65 identifies the represented lower constant with the actual activated
Gram and separates label erasure.  TPC--67 gives only a conditional
determinant route to the second arrow.  TPC--68 provides exact abstract
certificates for the fourth arrow but no physical arithmetic premise.
TPC--69--TPC--71 analyze only the first mixed trace, and
\cref{prop:four-corners-not-entry} prevents that scalar analysis from being
promoted to the fourth arrow.
The mass-degree comparison in TPC--68 requires the third arrow, while
\eqref{eq:tpc72-shorted-gram} shows that the fifth is independent.  None of
these papers proves the downstream boundary and tail estimates.  Finally,
TPC--63--TPC--65 record the strict endpoint budget.  Each item is therefore
a distinct necessary input in the declared staged architecture.
\end{proof}

\subsection{Non-substitution laws}

\begin{corollary}[No cross-branch substitution]
\label{cor:no-cross-branch-substitution}
The following implications are invalid without additional hypotheses:
\begin{align*}
 \text{determinant mass}
 &\not\Longrightarrow
 \text{missing-native operator bound},\\
 \text{primitive-plane scalar cancellation}
 &\not\Longrightarrow
 \text{represented lower frame},\\
 \text{anchored primitive-plane bound}
 &\not\Longrightarrow
 \text{entrywise erasure bound},\\
 \|R+E\|=o(1)
 &\not\Longrightarrow
 \|C_N\|=o(1),\\
 \tr(ER)=o(1)
 &\not\Longrightarrow
 \text{uniform row-star saving}.
\end{align*}
\end{corollary}

\begin{proof}
The first two statements concern different maps.  The third is
\cref{prop:four-corners-not-entry}.  The fourth follows because
\(C_N\) depends on \(B(N_A)\), which is not determined by the canonical
representative of \(M\) on \((\Ker A)^\perp\).  The fifth is disproved by
orthogonal direct sums of reinforcing and cancelling local stars, as in the
TPC--69 countermodels \citep{WangTPC69}.
\end{proof}
